\vspace{20pt}

Question: \textit{For the top-performing EPS sector, which individual companies drove that improvement — and were any consistent outliers?}

\vspace{20pt}

The analytics agent completed \textbf{three reasoning--execution cycles} before producing its final answer. The first cycle allowed the agent to understand the dataset schema for all five datasets. It loaded the datasets and determined which columns were important for analysis. The second cycle focused on exploratory analysis of Utilities-sector earnings-per-share (EPS) performance, while the third cycle was triggered by a validation step that identified inconsistencies in the initial ranking methodology and required a more rigorous reconciliation.

\subsect{Exploratory Per-Ticker EPS Analysis}

The agent began by reasoning about how to identify which Utilities companies contributed most to the sector's EPS improvement between 2014 and 2018. It proposed constructing a longitudinal dataset by concatenating the annual financial files and adding a \texttt{Year} column to each record. After filtering for the Utilities sector, it planned to compute several complementary metrics for every ticker:

\begin{itemize}
  \item Total EPS change between the first and last observed years,
  \item EPS compound annual growth rate (CAGR),
  \item Mean and standard deviation of reported EPS growth,
  \item Number of years with positive EPS growth.
\end{itemize}

The corresponding Python execution performed the following actions:

\begin{enumerate}
  \item Loaded and concatenated all five annual datasets (2014--2018).
  \item Standardized column names and converted EPS-related fields to numeric values.
  \item Filtered records where \texttt{Sector = Utilities}.
  \item Aggregated observations by ticker.
  \item Computed total EPS change, average EPS growth, growth volatility, and growth consistency metrics.
  \item Ranked companies by total EPS change and average EPS growth.
  \item Identified consistent positive, consistent negative, and highly volatile outliers.
\end{enumerate}

The execution produced metrics for approximately 110 Utilities companies. From these results, the agent observed that several firms exhibited very large endpoint EPS improvements and that some companies maintained positive EPS growth across all observed years. It also detected a set of volatility outliers based on unusually large standard deviations in EPS growth.

These observations directly influenced the next reasoning step. Rather than performing additional computation, the agent concluded that sufficient evidence existed to summarize top contributors, consistent positive performers, and volatile outliers. Consequently, the third step of the first cycle consisted only of narrative synthesis and no further code execution.

\subsect{Validation Feedback and Problem Identification}

After the first answer was generated, an independent validation stage compared the results against an alternative execution path. The validator detected a significant discrepancy in the list of top Utilities contributors. Specifically, the two analyses disagreed on the firms with the largest EPS improvements.

The validator concluded that the disagreement stemmed from inconsistent definitions of the ranking metric and instructed the agent to:

\begin{itemize}
  \item Use a single, fixed endpoint definition for EPS change,
  \item Require consistent inclusion criteria across all companies,
  \item Produce a ranked table with transparent calculations,
  \item Recompute consistency and volatility classifications,
  \item Document any data-quality issues.
\end{itemize}

This observation fundamentally altered the subsequent reasoning process. Rather than refining the previous summary, the agent decided to reconstruct the analysis from first principles using stricter definitions.

\subsect{Reconciliation and Robust Ranking}

The second reasoning cycle began with a revised analytical plan designed to eliminate ambiguity. The agent explicitly defined:

\[
  \text{Total\_EPS\_Change} = \text{EPS}_{2018} - \text{EPS}_{2014}
\]

and required that every ranked ticker possess valid EPS observations in both endpoint years.

The second Python execution carried out the following actions:

\begin{enumerate}
  \item Loaded all five annual datasets separately.
  \item Constructed a ticker-level panel containing EPS values for each year.
  \item Determined a representative sector label for every ticker.
  \item Filtered to Utilities companies only.
  \item Excluded firms missing either 2014 or 2018 EPS values.
  \item Computed endpoint EPS change using the fixed definition.
  \item Calculated the number of years with reported EPS.
  \item Flagged firms with large negative starting EPS values (\texttt{EPS\_2014 $\leq$ -1.0}).
  \item Computed EPS volatility using the standard deviation of yearly EPS values.
  \item Converted volatility values into z-scores and flagged outliers with $z > 2$.
  \item Identified firms with positive EPS in all five years.
  \item Generated a top-20 ranking table sorted by endpoint EPS improvement.
\end{enumerate}

The resulting observations substantially changed the interpretation of sector performance. The agent found that the largest numerical EPS improvements were frequently driven by companies recovering from extremely negative starting EPS values. For example, firms such as OPTT, EDN, and FCEL ranked highly by total EPS change but were flagged as recovery-driven cases because their 2014 EPS values were deeply negative.

At the same time, the analysis identified companies such as NEE and PAM that exhibited positive EPS throughout the entire 2014--2018 period while also producing meaningful EPS improvements. These firms were interpreted as more credible indicators of sustained sector strength.

The volatility analysis further influenced the final narrative. Companies such as CEQP, OPTT, and VST exhibited EPS volatility more than two standard deviations above the Utilities-sector average and were therefore classified as unstable outliers rather than representative sector drivers.

\subsect{Final Reasoning Outcome}

The final answer was generated without additional execution because the second computation resolved the inconsistency identified during validation. The agent's reasoning evolved from an initial exploratory ranking based on multiple growth metrics to a more rigorous endpoint-based attribution framework.

Overall, the execution trace demonstrates a two-stage analytical workflow:

\begin{enumerate}
  \item \textbf{Exploration and hypothesis generation}: identify candidate contributors and outliers using growth-based metrics.
  \item \textbf{Validation-driven reconciliation}: enforce a consistent definition of EPS improvement, perform data-quality checks, and distinguish genuine long-term contributors from recovery-driven or highly volatile firms.
\end{enumerate}

This iterative process allowed the agent to refine its interpretation and produce a final conclusion that emphasized sustained EPS improvement rather than merely large numerical recoveries.
